\section*{Chapter \ref{ch:ml:linear-and-regularised-baselines} --- Linear and Regularised Baselines}

\begin{solution}{exo:ml:linear-and-regularised-baselines:1}
Each side sells 30 positions of $1/50$ and buys 30: $0.6 + 0.6 = 1.2$ traded per side, 2.4 in all per unit of capital,
which is 1.2 times the gross book of 2. Cost: $2.4\times20\,\text{bp} = 0.48\%$ of capital that month.
\end{solution}

\begin{solution}{exo:ml:linear-and-regularised-baselines:2}
$1.6/2.5\times\sqrt{12} = 2.22$; after costs $1.4/2.5\times\sqrt{12} = 1.94$.
\end{solution}

\begin{solution}{exo:ml:linear-and-regularised-baselines:3}
Along $(1, 1)/\sqrt2$, the common part, which carries $(1 + 0.9)/2 = 95\%$ of the variance. It is the right direction
for a regression only if the target loads on the common part; if it loads on the difference of the two features, PCA's
last component is the one that matters and PCR with one component finds nothing.
\end{solution}

\begin{solution}{exo:ml:linear-and-regularised-baselines:4}
A decile portfolio uses only the extremes of the ranking, where the two models mostly agree, and it ignores the size of
the forecasts, where the trees' better R-squared partly lives. The trees' forecasts also move more (turnover 0.97
against 0.80), so they pay more costs: gross 2.20 against 2.17, net 1.90 against 1.91.
\end{solution}

\begin{solution}{exo:ml:linear-and-regularised-baselines:5}
The decile portfolio depends on the ranking (the IC); R-squared also depends on the scale of the forecasts. With two
years of data the trees rank tolerably but their forecasts spread too widely, more than twice the optimal scale
$c^\star$ of \cref{prop:ml:why:ceiling}, so R-squared is negative while the ranking earns a Sharpe ratio of 1.11.
\end{solution}

\begin{solution}{exo:ml:linear-and-regularised-baselines:6}
A comparison is only valid on the same data, splits, target, benchmark (zero or the mean) and period: R-squared varies
by several tenths of a point from one decade to the next (chapter~1). Refit the linear baseline in your own harness and
compare there, with the standard error of the monthly difference.
\end{solution}

\begin{solution}{exo:ml:linear-and-regularised-baselines:7}
\texttt{table(style=0.0)}: net Sharpe ratios 3.68 for ridge and 4.86 for the trees (4.24 for the lasso), against 1.59
and 1.90. R-squared moves less (0.35\% and 0.57\% against 0.30\% and 0.47\%). The factor returns add variance that a
decile portfolio, sorted on the characteristics that carry them, cannot diversify; R-squared counts every stock's
variance, of which factor risk is a small part.
\end{solution}

\begin{solution}{exo:ml:linear-and-regularised-baselines:8}
With $X^\top X = I$ the lasso objective separates: $\frac12(b_j - \beta_j)^2 + \lambda|\beta_j|$ for each least-squares
coefficient $b_j$, minimised by $\sign(b_j)\max(|b_j| - \lambda, 0)$. Once $\lambda$ exceeds the largest $|b_j|$, every
coefficient is zero and the model predicts the (zero) mean of the demeaned target, whose R-squared against zero is zero: the path's value at
$\lambda = 3\times10^{-3}$. Ranked characteristics are close to orthogonal, so the same logic holds approximately.
\end{solution}

\begin{solution}{pb:ml:linear-and-regularised-baselines:1}
\begin{enumerate}
\item 0.71\% in the test years; characteristics 0 to 5.
\item So that portfolios sorted on characteristics carry factor risk, as real ones do; they bring decile Sharpe ratios
from near 5 down to about 2.
\item 2.45 for their best network, 0.83 for a three-variable linear model.
\item The factor returns make the months, not the rows, the effective sample (chapter~1).
\item Ridge $\lambda = 10^4$; lasso $10^{-3}$ (six coefficients kept); at $3\times10^{-3}$ the lasso keeps none and scores
zero.
\item Ranked characteristics are nearly uncorrelated with equal variance, so principal components are arbitrary
directions (\cref{prop:ml:baselines:pls}); the signal is spread over many of them.
\item Its first component is $X^\top y$, the least-squares direction up to scale when the features are nearly orthogonal.
\item The square and the absolute value (additive nonlinearities), not the interaction $c_0c_3$.
\item Trees 0.47\%, IC 0.076, net Sharpe 1.90; lasso 0.32\%, IC 0.066, net Sharpe 1.91.
\item R-squared values sizes and all stocks; the decile portfolio values the extremes of the ranking, net of costs, and the
trees trade more.
\item Between five and ten years (60 months: $-0.02\%$ against 0.18\%; 120 months: 0.41\% against 0.34\%).
\item R-squared $-0.71\%$ but a net Sharpe ratio of 1.11 (ridge 0.90): the ranking is useful, the sizes are not.
\item \emph{Named result}: ridge 0.30\% and IC 0.060 against the trees' 0.47\% and 0.076 (the lasso: 0.32\% and 0.066);
the trees overtake ridge in R-squared between five and ten years of history, and never beat the lasso's net Sharpe
ratio.
\item The lasso: equal net performance, lower turnover, six coefficients that can be read and monitored.
\item Ridge on many random nonlinear features of the characteristics, heavily penalised.
\item The monthly net difference has a mean of $-0.08\%$ and a standard deviation of 2.56\%: $(2\times2.56/0.08)^2$ is
over 4\,000 months; no feasible track separates them.
\item Characteristics correlated in blocks, with the signal on the common part of each block.
\item The net Sharpe ratio and the turnover, beside the IC with its $t$.
\item The size of each move, which is mostly noise here: the \omterm{def:ml:linear-and-regularised-baselines:logistic}{logistic regression} ranks almost as well (IC 0.060).
\item A baseline prices complexity: it is the number a new model must beat, on equal terms, by more than the noise.
\end{enumerate}
\end{solution}

\begin{solution}{iq:ml:linear-and-regularised-baselines:1}
A penalised linear regression (ridge or lasso) on features ranked by date, with the target demeaned by date, the penalty
chosen by purged cross-validation on whole dates, and the standard report (R-squared against zero, IC with HAC $t$,
decile portfolio net of costs). Every later model is judged against it.
\iqlookfor{ranking, demeaning, purged folds on dates, and a report beyond R-squared.}
\end{solution}

\begin{solution}{iq:ml:linear-and-regularised-baselines:2}
Try both: ridge if many characteristics carry small effects, lasso if a few carry most; with ranked, nearly orthogonal
features the lasso's sparsity is cheap and its model readable. The elastic net covers correlated groups. Choose on
purged cross-validation, not taste.
\iqlookfor{the sparsity prior, correlated features, and an empirical choice.}
\end{solution}

\begin{solution}{iq:ml:linear-and-regularised-baselines:3}
PCR keeps the directions of largest feature variance, ignoring the target; PLS keeps directions of largest covariance with
the target. They differ when the signal lies in low-variance directions or when features have similar variance (ranks),
where PCA's directions are arbitrary.
\iqlookfor{supervised versus unsupervised reduction and the case where it matters.}
\end{solution}

\begin{solution}{iq:ml:linear-and-regularised-baselines:4}
R-squared rewards correct sizes and every stock; a portfolio rewards the ranking at the extremes and pays for turnover.
A model can improve sizes in the middle of the distribution, or trade more, and leave the portfolio unchanged or worse.
\iqlookfor{the gap between forecast accuracy and trading value.}
\end{solution}

\begin{solution}{iq:ml:linear-and-regularised-baselines:5}
Rank features by date, demean the target by date, make purged folds of whole dates with an embargo of the label length,
evaluate a logarithmic grid wide enough that both ends are worse, and take the most penalised value within one standard
error of the best.
\iqlookfor{folds by date, a grid that brackets the optimum, and the one-standard-error rule.}
\end{solution}

\begin{solution}{iq:ml:linear-and-regularised-baselines:6}
Maximise $w^\top X^\top y$ over unit $w$: by Cauchy--Schwarz, $w\propto X^\top y$. The first principal component maximises
$w^\top X^\top Xw$ instead, which ignores $y$; when the target depends on a low-variance direction, PLS finds it and PCA
does not.
\iqlookfor{the derivation and the geometric reason.}
\end{solution}
